% SHARD-ID: AQM-83-QSA-PRODUCT-RINGS-TENSOR-FACTORS
% SHARD-TITLE: Product Rings and Tensor Factors
% SHARD-SUMMARY: Reviews product-field spectra and mixed-residue product warnings for the QSA catalogue.
% SHARD-SUMMARY: Compares product-ring residue-site tensor factors with the structureless finite-set tensor-site workflow.
% SHARD-SUMMARY: Records cotangent-first edge cases and database/helper gaps without claiming a universal product-ring equivalence.
% SHARD-KEYWORDS: quantum-system association, product ring, tensor factors, residue fields, cotangent, finite-ring database

\section{Product Rings and Tensor Factors}
\label{sec:qsa-product-rings-tensor-factors}

\subsection{Status and Local Anchors}

\textbf{Status: Review shard.}  This is QSA Step 19 from
\path{docs/quantum_system_associations/PLAN.md}.  It adds no source,
convention entry, script, run bundle, generated data, populated database row,
or finite-ring decomposition theorem.  It reviews exactly the product-ring
evidence already in AQM-40, AQM-58, AQM-68, AQM-80, AQM-81, and AQM-82.

The governing conventions are \texttt{(ad)} for product-field spectra,
\texttt{(an)} for finite-ring database and helper scope, \texttt{(ar)} for
structureless finite-set tensor sites, and \texttt{(au)} for QSA table status
tokens.  The product-field anchor is
\[
  R=k^n,\qquad k=\mathbb F_q,
\]
where AQM-40 proves that \(\Spec(k^n)\) is the finite discrete \(n\)-point
spectrum with residue field \(k\) at every point.  The main comparison anchor
is AQM-58: the same product ring has
\[
  \Omega^1_{k^n/k}=0,
\]
so the cotangent-first tangent-Weyl layer can vanish even when the residue
workflow has \(n\) qudit tensor factors.

\subsection{Selected Workflow: Product-Field Residue Sites}

In the selected product-field residue-site workflow, the ring supplies more
than a bare set but less than a universal quantum model.  For
\[
  R=k^n
\]
the idempotents and projections select the points
\[
  x_i=\ker(\operatorname{pr}_i),\qquad i=1,\ldots,n,
\]
with \(\kappa(x_i)=k\).  Convention \texttt{(ad)} and AQM-40 then give
\[
  E(X)=\bigoplus_{i=1}^n k^2,\qquad
  \mathcal H(X)=\bigotimes_{i=1}^n\ell^2(k),\qquad
  \mathcal A(X)=\bigotimes_{i=1}^n
    \operatorname{End}_{\mathbb C}\ell^2(k).
\]
Thus \(\mathbb F_3^n\) gives \(n\) qutrit residue factors, not
\(3^n\) spectrum points.  Open subsets are just subsets of the \(n\)
points, and the local observable inclusions are tensoring by identities on
the missing residue factors.

This is a tensor factorization statement for this residue-site association
only.  It does not identify product rings with all first-association,
cotangent-first, field-label, nilpotent-sensitive, stabilizer, Clifford, or
dynamics workflows.  AQM-40 already separates the ambient \(n\)-qudit Weyl
net from the later choice of an isotropic stabilizer subgroup and phases.

\subsection{Comparison With Structureless Tensor Sites}

AQM-68 starts with a bare finite set \(X\) and only then adds site data
\[
  \{\mathcal H_x\}_{x\in X}
  \quad\text{or}\quad
  \{\mathcal A_x\}_{x\in X}.
\]
After those choices, it forms
\[
  \mathcal H(S)=\bigotimes_{x\in S}\mathcal H_x,\qquad
  \mathcal A(S)=\bigotimes_{x\in S}\mathcal A_x.
\]

The product-field workflow supplies the following extra data that the bare
finite set did not supply.

\begin{itemize}
\item The site set is not arbitrary: it is the spectrum
  \(\{x_i=\ker(\operatorname{pr}_i)\}\) of the displayed product ring.
\item The topology is discrete for \(k^n\), with singleton opens \(D(e_i)\).
\item Each site has a residue field \(\kappa(x_i)=k\), hence a local
  finite-field phase space \(k^2\).
\item The local carrier is the particular Weyl carrier
  \(\ell^2(\kappa(x_i))\), and the local algebra is its full endomorphism
  algebra.
\end{itemize}

The analogy with AQM-68 is therefore precise but limited: both workflows use
finite tensor products and tensor-by-identity inclusions once the local
factors are known.  They do not have the same input.  The structureless
finite set allows arbitrary chosen Hilbert spaces; the product-field
residue-site workflow fixes the local Weyl factors from the residue fields.

\subsection{Mixed Residue Warning}

AQM-23 and AQM-81 supply the warning example
\[
  R=\mathbb Z/6\mathbb Z.
\]
The residue-site workflow has two points, with residues
\[
  \mathbb F_2,\qquad \mathbb F_3,
\]
and hence
\[
  E=\mathbb F_2^2\oplus\mathbb F_3^2,\qquad
  \mathcal H=\ell^2(\mathbb F_2)\otimes\ell^2(\mathbb F_3).
\]
This is still an independent tensor-factor residue workflow, but it is not
one common \(\mathbb F_3\)-linear or \(\mathbb F_2\)-linear phase space.  Any
common-prime-field reduction, stabilizer comparison, or Clifford comparison
would require its own pairing and phase convention.

This shard does not promote \(\mathbb Z/6\mathbb Z\) to a general
quotient-collapse theorem.  AQM-80 keeps quotient-collapse examples beyond
the existing \(\mathbb Z/6\mathbb Z\) anchor blocked by
\texttt{aqm-qsa-collapse01}.

\subsection{Cotangent Edge Cases}

Product-ring tensor factors are highly workflow-dependent.  The clearest
edge case is AQM-58.  For \(k^n\), every \(k\)-derivation out of \(k^n\)
vanishes, so
\[
  \Omega^1_{k^n/k}=0.
\]
In the tangent-cotangent Weyl workflow of conventions \texttt{(ao)} and
\texttt{(av)}, the local tangent and cotangent labels at each product point
therefore vanish.  The cotangent-first geometric Weyl carrier for those
directions is the zero-label, one-dimensional case; it is not the
\(n\)-qudit residue carrier of AQM-40.

The single field row already shows the same distinction in smaller form.
\(\Spec(\mathbb F_3)\) has one residue-site qutrit in AQM-81, while its
relative cotangent over \(\mathbb F_3\) is zero in AQM-77 and AQM-80.  The
product field repeats this zero-cotangent phenomenon \(n\) times while the
residue-site layer still has \(n\) local Weyl factors.

For mixed or nonreduced finite rings, the table in AQM-80 is the current
boundary.  \(\mathbb Z/6\mathbb Z\) has the mixed-residue tensor carrier
above, but no database-backed cotangent comparison row is imported here.
\(\mathbb Z/4\mathbb Z\) and dual-number examples have local cotangent or
nilpotent-sensitive anchors in AQM-77 and AQM-82, but they are not product
field tensor-factor examples in this Step 19 review.

\subsection{Database and Helper Boundary}

Convention \texttt{(an)} records a manual
\texttt{finite\_ring\_product} constructor for the finite-ring database
implementation: it preserves argument order, concatenates factor coordinate
blocks and identity vectors, copies factor multiplication into each block,
and sets cross-factor products to zero.  This is a data-layout convention for
MVP constructors, not a classification theorem for arbitrary finite rings.

The same convention limits the in-memory residue-quantisation helper to
source-backed cases already fixed in the lab book: prime fields, the
\(\mathbb Z/6\mathbb Z\) row from AQM-23, and explicit finite product fields
from AQM-40.  The SQLite smoke build remains schema-only, and the review CSVs
are in-memory MVP exports.  No populated or audited finite-ring database is
claimed.

The active gaps therefore remain:

\begin{description}
\item[\texttt{aqm-qsa-frres01}.]
Arbitrary finite-ring residue decomposition and site ordering.
\item[\texttt{aqm-qsa-frcot01}.]
Database-backed cotangent rows.
\item[\texttt{aqm-qsa-collapse01}.]
Quotient-collapse examples beyond the \(\mathbb Z/6\mathbb Z\) anchor.
\item[\texttt{aqm-qsa-artin01}.]
Unsupported thickened or multi-direction Artin layers.
\end{description}

In particular, a product of thickened local rings may have a displayed tensor
product after each local Artin/Frobenius layer has its own certified
generating character, but Step 19 does not assert such a general policy.

\subsection{Conclusion}

Product rings give independent tensor factors only in a selected and
well-anchored sense.  For \(k^n\), the product-field residue-site association
supplies \(n\) finite Weyl sites, \(n\) local matrix algebras, and the usual
tensor-by-identity local net.  This matches the form of the structureless
finite-set tensor-site construction only after the product ring has supplied
the sites and residue-field Weyl factors.

The same ring may look degenerate in a different association workflow:
cotangent-first labels vanish for \(k^n\), stabilizers and dynamics remain
extra choices, mixed residues are not a single common-field phase space, and
nilpotent-sensitive product layers require separate Artin/Frobenius evidence.
Those boundaries are part of the QSA catalogue rather than exceptions to a
universal product-ring tensor-factor rule.
